\documentclass[11pt]{article}
\usepackage[margin=1in]{geometry}
\usepackage{amsmath,amssymb,amsthm,lmodern}
\usepackage[colorlinks=true,linkcolor=blue,urlcolor=blue]{hyperref}
\newtheorem{theorem}{Theorem}
\setlength{\emergencystretch}{2em}
\title{An endpoint correction to a printed positivity-rigidity conjecture}
\author{Myeongjin Kang}
\date{25 October 2026}
\begin{document}
\maketitle
\begin{abstract}
The finite-sum rigidity statement in Conjecture 11 of Lóczi and Ketcheson constrains a coefficient that is absent from every hypothesis. We give a one-parameter family satisfying every premise and violating that endpoint conclusion, with a complete Lean certificate. This is a literal endpoint correction; it does not refute the intended rigidity statement restricted to the coefficients that occur in the sums.
\end{abstract}
\section{The printed assertion}
For $s\geq1$, nonzero real $q$, positive real $\widetilde q$, and real coefficients $p_0,\ldots,p_{s+1}$, define
\[
 S_k=\sum_{m=0}^{\min(k,s)}\binom{s-1+k-m}{s-1}p_mq^m(-\widetilde q)^{k-m}.
\]
Conjecture 11 in \cite{source}, recorded as MathDB 325243, assumes $S_k\geq0$ for $0\leq k\leq s+1$. It concludes $p_0\geq0$ and
\[
 p_m=p_0\binom{s}{m}(\widetilde q/q)^m\quad(1\leq m\leq s+1),
\]
and also $S_0=p_0$ and $S_k=0$ for $1\leq k\leq s+1$.
The upper coefficient index and the summation cutoff above agree with the original versioned source, not only with the database transcription.
\section{Complete counterexample}
\begin{theorem}
For every nonzero real $t$, the choices
\[
 s=1,\quad q=\widetilde q=1,\quad p_0=p_1=1,\quad p_2=t
\]
satisfy all printed assumptions and violate the printed conclusion.
\end{theorem}
\begin{proof}
The parameter restrictions hold. Since every binomial coefficient in the three sums is $\binom{k-m}{0}=1$, direct evaluation gives
\[
 S_0=1,\qquad S_1=-1+1=0,\qquad S_2=1-1=0.
\]
These are exactly all required inequalities, so every hypothesis holds. The claimed conclusion at $m=2$ would be
\[
 t=p_2=p_0\binom12(\widetilde q/q)^2=0,
\]
contrary to the choice of $t$.
\end{proof}
More generally, each summand has $m\leq s$. Hence changing $p_{s+1}$ never changes any $S_k$, for any $k$. The conclusion at that coefficient cannot follow without an additional condition.
\section{Scope and formal certificate}
The surrounding application of \cite{source} takes $p_m$ from derivatives of a polynomial whose degree is at most $s$, so its $(s+1)$st coefficient vanishes. The displayed free-coefficient conjecture does not impose this condition. It is plausible that the intended conclusion should stop at $m=s$, or that the endpoint coefficient was meant to be fixed at zero. This note makes no claim about the validity of either corrected formulation and is not a new substantive resolution of the intended rigidity problem. Its outcome is the exact correction to the printed statement.

The Lean file \nolinkurl{formal/PositivityRigidity.lean} defines the exact sum and all hypotheses and conclusions. A total coefficient sequence merely supplies values outside the finite range; those values are never used. The theorem \nolinkurl{sum_ignores_endpoint} proves the free-endpoint observation for arbitrary parameters. The remaining theorems prove every required inequality for the family above and refute the universal printed assertion in \nolinkurl{printed_conjecture_false}.

All eight local declarations compile with Lean 4.34.1 against pinned mathlib dependencies. The extracted release contains the independent review, dependency pins and \nolinkurl{verify_outcome.py}. Only the usual foundations \nolinkurl{propext}, \nolinkurl{Classical.choice}, and \nolinkurl{Quot.sound} are permitted. The result does not depend on a numerical search, an admitted lemma, or assumptions imported from the conjecture.
\begin{thebibliography}{9}
\bibitem{source} Lajos Lóczi and David I. Ketcheson. \emph{Rational functions with maximal radius of absolute monotonicity}. arXiv:1303.6651v1, 2013. Section 2.5, Conjecture 11. \url{https://arxiv.org/html/1303.6651v1}.
\end{thebibliography}
\end{document}
